% 02_part1_intro.tex — ส่วนที่ 1: ส่วนนำ

% ─── 1.1 บทสรุปผู้บริหาร ─────────────────────────────────────────────────────
\section{บทสรุปผู้บริหาร} \label{sec:exec_summary}

รายงานการประเมินตนเองการประกันคุณภาพการศึกษาภายใน ระดับหลักสูตร ตามเกณฑ์คุณภาพ AUN-QA ของหลักสูตรวิศวกรรมศาสตรมหาบัณฑิต สาขาวิชาวิศวกรรมคอมพิวเตอร์และปัญญาประดิษฐ์ คณะเทคโนโลยีอุตสาหกรรม มหาวิทยาลัยราชภัฏอุตรดิตถ์ ปีการศึกษา 2568 ซึ่งเป็นหลักสูตรใหม่ พ.ศ. 2568 มีนักศึกษารุ่นแรกจำนวน 4 คน ทั้งนี้ อาจารย์ผู้รับผิดชอบหลักสูตรทั้ง 3 คน สำเร็จการศึกษาระดับปริญญาเอกทั้งหมด (จำนวน 3 คน) และดำรงตำแหน่งทางวิชาการระดับผู้ช่วยศาสตราจารย์ทั้งหมด (จำนวน 3 คน) โดยผลการประเมินตนเองในภาพรวมของทั้ง 8 เกณฑ์ (Criteria) อยู่ในระดับคะแนน 2 ซึ่งแสดงถึงการดำเนินงานของหลักสูตรที่มีระบบและกลไกพื้นฐานรองรับการจัดการเรียนการสอนอย่างมีหลักฐานเชิงประจักษ์ตามเกณฑ์มาตรฐาน และมีความพร้อมสำหรับการพัฒนาคุณภาพการจัดการศึกษาอย่างต่อเนื่องตามรายละเอียดดังต่อไปนี้
\begin{table}[H]
  \centering
  \caption{ตารางการประเมินตนเองของหลักสูตร}
  \label{tbl:exec_summary}
  \begingroup
  \footnotesize
  \renewcommand{\arraystretch}{1.1}
  \begin{tabularx}{0.8\linewidth}{|l|>{\raggedright\arraybackslash}X|c|}
    \hline
    \rowcolor{gray!15}
    \multicolumn{2}{|c|}{\textbf{ตัวบ่งชี้ / Criteria}} & \textbf{ประเมินตนเอง} \\
    \hline
    Criterion 1 & Expected Learning Outcome & 2 \\
    \hline
    Criterion 2 & Programme Structure and Content & 2 \\
    \hline
    Criterion 3 & Teaching and Learning Approach & 2 \\
    \hline
    Criterion 4 & Student Assessment & 2 \\
    \hline
    Criterion 5 & Academic Staff & 2 \\
    \hline
    Criterion 6 & Student Support Service & 2 \\
    \hline
    Criterion 7 & Facilities and Infrastructure & 2 \\
    \hline
    Criterion 8 & Output and Outcomes & 2 \\
    \hline
    \rowcolor{gray!10}
    \multicolumn{2}{|c|}{\textbf{Overall}} & \textbf{2} \\
    \hline
  \end{tabularx}
  \endgroup
  \par\vspace{0.3em}
  \noindent{\footnotesize \textbf{หมายเหตุ}: รายละเอียดและเหตุผลประกอบการประเมินตนเองของแต่ละเกณฑ์ ปรากฏในตารางสรุปผลการประเมินตนเอง ส่วนที่ 3 หน้า~\pageref{tbl:sar_summary}}
\end{table}

% ─── 1.2 ข้อมูลพื้นฐาน ───────────────────────────────────────────────────────
\section{ข้อมูลพื้นฐาน} \label{sec:basic_info}

\subsection{ข้อมูลพื้นฐานของมหาวิทยาลัย}

\subsubsection*{ปรัชญามหาวิทยาลัย}

``มหาวิทยาลัยเพื่อการศึกษาและพัฒนาชุมชนท้องถิ่น''

\subsubsection*{ปรัชญาการจัดการศึกษา}

``สร้างผลลัพธ์การเรียนรู้ด้วยประสบการณ์เชิงบูรณาการ'' ``Integrated Experiential Learning (IEL)''

\textbf{``ผลลัพธ์การเรียนรู้'' (Learning Outcomes)} หมายถึง ผลที่เกิดขึ้นแก่ผู้เรียนผ่านกระบวนการเรียนรู้
ที่ได้จากการศึกษาฝึกอบรม หรือประสบการณ์ที่เกิดขึ้นจากการฝึกปฏิบัติหรือการเรียนรู้จริงในที่ทำงานระหว่าง
การศึกษา ประกอบด้วยผลลัพธ์ 4 ด้าน ได้แก่ ด้านความรู้ (Knowledge) ด้านทักษะ (Skills)
ด้านจริยธรรม (Ethics) และด้านลักษณะบุคคล (Character)

\textbf{``ประสบการณ์เชิงบูรณาการ'' (Integrated Field Experiences)} หมายถึง กระบวนการเรียนรู้ที่เกิด
จากการเรียนรู้และบูรณาการทางวิชาการ วิชาชีพ การฝึกปฏิบัติระหว่างการศึกษาและการฝึกประสบการณ์
เพื่อพัฒนาตนเอง พัฒนาท้องถิ่น สู่การเรียนรู้ตลอดชีวิต

\textbf{``ผลลัพธ์การเรียนรู้เฉพาะด้าน'' (Specific Learning Outcomes)} หมายถึง ผลที่เกิดกับผู้เรียนผ่าน
กระบวนการเรียนรู้รายวิชาเฉพาะด้านของแต่ละหลักสูตร กระบวนการศึกษา การฝึกอบรม การฝึกปฏิบัติหรือ
การเรียนรู้จริงระหว่างการศึกษา การฝึกเป็นนักคิด นักจัดการ รวมถึงการฝึกประสบการณ์ในสถานประกอบการ

\textbf{``ผลลัพธ์การเรียนรู้ทั่วไป'' (Generic Learning Outcomes)} หมายถึง มุ่งพัฒนาผู้เรียนให้เป็นผู้ใฝ่เรียน
แสวงหาความรู้เพื่อพัฒนาตนเอง สร้างความรู้คู่คุณธรรม เคารพคุณค่าและศักดิ์ศรีความเป็นมนุษย์
มีความรับผิดชอบต่อสังคมและสิ่งแวดล้อม เป็นพลเมืองดิจิทัลและพลเมืองเข้มแข็ง

\subsubsection*{วิสัยทัศน์}

``มหาวิทยาลัยพันธกิจสัมพันธ์ที่มีคุณภาพ สร้างคุณค่าเพื่อพัฒนาท้องถิ่น''

\subsubsection*{อัตลักษณ์บัณฑิต}

``บัณฑิตนักคิด นักจัดการ สร้างงานด้วยความรู้คู่คุณธรรม''

\subsubsection*{เอกลักษณ์มหาวิทยาลัย}

``มหาวิทยาลัยราชภัฏอุตรดิตถ์เป็นปัญญาของชุมชน'' (\textit{Wisdom of Community})

\subsubsection*{ค่านิยมองค์กร (URU Core Values)}

\begin{itemize}
  \item \textbf{U --- Unity} ความเป็นหนึ่งเดียว
  \item \textbf{R --- Responsibility} ความรับผิดชอบ
  \item \textbf{U --- Uniqueness} ความเป็นเอกลักษณ์มหาวิทยาลัยพันธกิจสัมพันธ์
\end{itemize}

\subsubsection*{พันธกิจหลัก}

\begin{enumerate}
  \item ผลิตบัณฑิตดีที่มีคุณภาพ มีทัศนคติที่ดีเป็นพลเมืองดีในสังคมและมีสมรรถนะตามความต้องการของผู้ใช้บัณฑิต
  \item ผลิตและพัฒนาครูอย่างมีคุณภาพมาตรฐานของคุรุสภา
  \item วิจัยและบริการวิชาการ ถ่ายทอดเทคโนโลยี สร้างองค์ความรู้และนวัตกรรมที่มีคุณภาพและได้มาตรฐานเป็นที่ยอมรับ มุ่งเน้นการบูรณาการเพื่อนำไปใช้ประโยชน์ในท้องถิ่นได้อย่างเป็นรูปธรรม
  \item พัฒนาท้องถิ่นและทำนุบำรุงศิลปวัฒนธรรมตามศักยภาพ สภาพปัญหาและความต้องการที่แท้จริงของชุมชน โดยน้อมนำแนวพระราชดำริสู่การปฏิบัติ
  \item สร้างเครือข่ายความร่วมมือกับทุกภาคส่วนเพื่อพัฒนาท้องถิ่นและเสริมสร้างความเข้มแข็งของผู้นำชุมชนให้มีคุณภาพและความสามารถในการบริหารงานเพื่อประโยชน์ต่อชุมชนท้องถิ่น
  \item พัฒนาระบบบริหารจัดการมหาวิทยาลัยด้วยหลักธรรมาภิบาลพร้อมรับการเปลี่ยนแปลง เพื่อให้เกิดการพัฒนาต่อเนื่อง
\end{enumerate}

\subsubsection*{ยุทธศาสตร์มหาวิทยาลัยราชภัฏอุตรดิตถ์ ระยะ 5 ปี (พ.ศ. 2565--2569)}

\textbf{ยุทธศาสตร์ที่ 1 ขับเคลื่อนการดำเนินงานตามแผนแม่บท}\\
มีเป้าประสงค์หลัก คือ การดำเนินการพัฒนาท้องถิ่นในด้านเศรษฐกิจ สังคม สิ่งแวดล้อม และการศึกษา

\textbf{ยุทธศาสตร์ที่ 2 การพัฒนาพันธกิจหลักเพื่อขับเคลื่อนสู่การเป็นมหาวิทยาลัยพันธกิจสัมพันธ์}\\
มีเป้าประสงค์หลัก คือ การผลิตบัณฑิตและการยกระดับคุณภาพการศึกษา การพัฒนาด้านหลักสูตรและการจัดการศึกษา ด้านการพัฒนากำลังคนเพื่อรองรับความเปลี่ยนแปลงของประเทศ ด้านการพัฒนางานวิจัย เทคโนโลยีและนวัตกรรม ด้านการบริการวิชาการและศิลปวัฒนธรรม และโครงการอันเนื่องมาจากพระราชดำริ

\textbf{ยุทธศาสตร์ที่ 3 การพัฒนาการบริหารจัดการสู่องค์กรที่มีสมรรถนะสูง}\\
มีเป้าประสงค์หลัก คือ การพัฒนาระบบบริหารจัดการสู่ความเป็นเลิศ ด้านการพัฒนากายภาพและด้านบุคลากร

\textbf{ยุทธศาสตร์ที่ 4 การพลิกโฉมมหาวิทยาลัยสู่ความเป็นเลิศ}\\
มีเป้าประสงค์หลัก คือ การสนับสนุนการจัดการศึกษาตลอดชีวิต การพัฒนาทักษะเพื่ออนาคตในรูปแบบ non-degree และ degree เพิ่มทักษะด้านการประกอบอาชีพและการทำงาน การพัฒนาคน สถาบันความรู้ชุมชนเชิงพื้นที่ มุ่งเน้นการสร้างและพัฒนาบุคลากรที่มีทักษะสูงตามความต้องการของท้องถิ่น

\subsection{ข้อมูลพื้นฐานของคณะ}

หลักสูตร วศ.ม. วิศวกรรมคอมพิวเตอร์และปัญญาประดิษฐ์ อยู่ภายใต้ความรับผิดชอบร่วมของ
2 คณะ ได้แก่ คณะเทคโนโลยีอุตสาหกรรม (คณะหลัก) และคณะวิทยาศาสตร์และเทคโนโลยี (คณะร่วม)

\subsubsection*{คณะเทคโนโลยีอุตสาหกรรม (คณะหลัก)}

\textbf{ปรัชญา}\\
``สร้างบัณฑิตสู่งาน เชี่ยวชาญเทคโนโลยี มีทักษะเป็นเลิศ เชิดชูคุณธรรม น้อมนำนวัตกรรม เพื่อส่งเสริมสังคม''

\textbf{ปณิธาน}\\
``สุจริตต่อตน มั่นคงต่อหน้าที่ ศึกษาเทคโนโลยีอย่างถ่องแท้ แก้ปัญหาอย่างมีคุณธรรมนำทาง''

\textbf{วิสัยทัศน์}\\
``คณะเทคโนโลยีอุตสาหกรรม เป็นองค์กรคุณภาพที่ผลิตบัณฑิตให้เป็นเลิศด้านเทคโนโลยี สนับสนุนการวิจัยสมัยใหม่ และเป็นศูนย์กลางการถ่ายทอดวิทยาการแก่ชุมชน''

\textbf{อัตลักษณ์บัณฑิต}\\
``บัณฑิตใฝ่วิชาการ อดทนสู้งาน นักปฏิบัติการ จิตใจคุณธรรมนำวิชาชีพ''

\vspace{6pt}
\subsubsection*{คณะวิทยาศาสตร์และเทคโนโลยี (คณะร่วม)}

\textbf{ปรัชญาคณะ}\\
``วิทยาศาสตร์และเทคโนโลยีเพื่อการพัฒนาท้องถิ่นและประเทศ ด้วยองค์ความรู้ นวัตกรรมและความเป็นหนึ่งเดียวกับชุมชน''

\textbf{วิสัยทัศน์คณะ}\\
``คณะวิทยาศาสตร์และเทคโนโลยีจะเป็นองค์กรแห่งความเป็นหนึ่งเดียว เป็นศูนย์กลางองค์ความรู้และนวัตกรรมทางวิทยาศาสตร์และเทคโนโลยี สร้างบัณฑิตคุณภาพที่มีคุณธรรม จิตอาสา ความรู้ และทักษะวิชาชีพ เพื่อเป็นพลังขับเคลื่อนการพัฒนาท้องถิ่นและชุมชนสู่ความยั่งยืน''

\textbf{ภารกิจหลักของคณะ}

\begin{enumerate}
  \item ผลิตบัณฑิตคุณภาพที่มีคุณธรรม จิตอาสา ความรู้และทักษะวิชาชีพ ตอบสนองความต้องการของสังคมและตลาดแรงงาน บนพื้นฐานปรัชญาเศรษฐกิจพอเพียง
  \item บูรณาการวิทยาศาสตร์และเทคโนโลยีเพื่อขับเคลื่อนเศรษฐกิจ BCG และเป้าหมายการพัฒนาที่ยั่งยืน มุ่งสู่การเป็นคณะวิทยาศาสตร์และเทคโนโลยีสีเขียว (Green Faculty) ที่มีการจัดการทรัพยากรอย่างมีประสิทธิภาพ
  \item สร้างองค์ความรู้ วิจัย และนวัตกรรมที่บูรณาการศาสตร์กับภาคีเครือข่าย เพื่อนำไปใช้ประโยชน์ในการพัฒนาท้องถิ่น ประเทศ และสังคมอย่างยั่งยืน
  \item ส่งเสริมกิจกรรมนักศึกษาและการเรียนรู้นอกห้องเรียน เพื่อปลูกฝังจิตอาสา ภาวะผู้นำ และทักษะการทำงานร่วมกันบนแนวคิดวิศวกรสังคม
  \item บริหารจัดการองค์กรด้วยหลักธรรมาภิบาล โปร่งใส ตรวจสอบได้ และปรับตัวรองรับการเปลี่ยนแปลง
  \item พัฒนาระบบเทคโนโลยีสารสนเทศและโครงสร้างพื้นฐาน เพื่อสนับสนุนการบริหาร การเรียนการสอน และการวิจัยอย่างมีประสิทธิภาพ
  \item สร้างเครือข่ายความร่วมมือกับชุมชน ภาคเอกชน และองค์กรต่างๆ เพื่อเพิ่มโอกาส สร้างคุณค่า และรายได้ที่เกื้อกูลต่อบุคลากร นักศึกษา และคณะ
  \item พัฒนากลยุทธ์การประชาสัมพันธ์และการรับนักศึกษา โดยใช้เครือข่าย ศิษย์เก่า และพันธมิตรทางวิชาการ เพื่อเข้าถึงกลุ่มเป้าหมายอย่างมีประสิทธิภาพ
\end{enumerate}

\textbf{อัตลักษณ์บัณฑิต}\\
``บัณฑิตคุณภาพ มีคุณธรรมและจิตอาสา เชี่ยวชาญวิทยาศาสตร์และเทคโนโลยี พร้อมทักษะวิชาชีพเพื่อพัฒนาท้องถิ่นอย่างยั่งยืน''

\textbf{ยุทธศาสตร์คณะวิทยาศาสตร์และเทคโนโลยี}

\textbf{ยุทธศาสตร์ที่ 1 การผลิตบัณฑิตคุณภาพ} --- พัฒนาทักษะภาษาและดิจิทัล; ทักษะวิชาการและวิชาชีพเพื่อตลาดแรงงาน; ทักษะวิจัยและนวัตกรรมร่วมกับชุมชน; ทักษะชีวิตและจิตอาสาเชิงวิศวกรสังคม; สร้างเครือข่ายศิษย์เก่าและพัฒนาศักยภาพต่อเนื่อง

\textbf{ยุทธศาสตร์ที่ 2 การพัฒนาหลักสูตรและการจัดการเรียนการสอน} --- พัฒนาหลักสูตรคุณภาพตาม OBE และ WIL/CWIE; หลักสูตรตอบสนองผู้เรียนและตลาดแรงงาน (รวม Life Long Learning / Credit Bank); พัฒนาห้องเรียน ห้องปฏิบัติการ และความปลอดภัย Lab

\textbf{ยุทธศาสตร์ที่ 3 การพัฒนาวิทยาศาสตร์และเทคโนโลยีที่ยั่งยืน} --- ขับเคลื่อนเศรษฐกิจ BCG ด้วยวิทยาศาสตร์และเทคโนโลยี; พัฒนาคณะสู่ Green Faculty และ Green University; บูรณาการงานวิจัยและบริการวิชาการกับ SDGs

\textbf{ยุทธศาสตร์ที่ 4 การวิจัยและบริการวิชาการเพื่อท้องถิ่นและประเทศ} --- ส่งเสริมงานวิจัย บริการวิชาการและพันธกิจสัมพันธ์; ยกระดับคุณภาพผลงานวิจัยและบริการวิชาการ; บูรณาการงานวิจัยกับการเรียนการสอนและการฝึกประสบการณ์

\textbf{ยุทธศาสตร์ที่ 5 การบริหารจัดการที่มีธรรมาภิบาล และการพัฒนาศักยภาพบุคลากร} --- พัฒนาระบบบริหารจัดการที่มีธรรมาภิบาลและวัฒนธรรมองค์กรแห่งความเป็นหนึ่งเดียว; ส่งเสริมศักยภาพบุคลากรสายวิชาการและสายสนับสนุนด้านวิชาการ; ส่งเสริมสุขภาวะและความปลอดภัยทางจิตใจของบุคลากร

\textbf{ยุทธศาสตร์ที่ 6 สืบสานประเพณี ทำนุบำรุงศิลปวัฒนธรรม ถ่ายทอดภูมิปัญญาท้องถิ่นและโครงการอันเนื่องมาจากพระราชดำริ} --- เชิดชูกิจกรรมสืบสานและทำนุบำรุง; ร่วมมือชุมชนด้านศิลปวัฒนธรรม; เก็บรวบรวมและถ่ายทอดภูมิปัญญาท้องถิ่น; เชื่อมโยงโครงการอันเนื่องมาจากพระราชดำริ

\textbf{ยุทธศาสตร์ที่ 7 การพัฒนาระบบเทคโนโลยีสารสนเทศและโครงสร้างพื้นฐาน} --- พัฒนาระบบสารสนเทศเพื่อการบริหารจัดการ (Dashboard คลังปัญญา Resume Builder ติดตามศิษย์เก่า บันทึกโครงการบริการวิชาการ); โครงสร้างพื้นฐานอัจฉริยะและมหาวิทยาลัยสีเขียว; ระบบช่างประจำคณะและงานซ่อมบำรุง

\textbf{ยุทธศาสตร์ที่ 8 การสร้างเครือข่ายความร่วมมือและการสร้างรายได้} --- สร้างเครือข่ายความร่วมมือเชิงยุทธศาสตร์ (MOU/MOA และโครงการร่วม); เพิ่มรายได้จากศักยภาพวิชาการ บุคลากร และโครงสร้างพื้นฐานของคณะ

\textbf{ยุทธศาสตร์ที่ 9 การประชาสัมพันธ์และการรับนักศึกษา} --- สร้างภาพลักษณ์คณะที่รู้จักและน่าเชื่อถือผ่านสื่อสมัยใหม่; บริหารจัดการและเครือข่ายรับนักศึกษา (โรงเรียนเครือข่าย แผนปฏิบัติงาน กิจกรรมแนะแนว การเล่าเรื่องศิษย์เก่า)

% ─── 1.2.3 หลักสูตร ──────────────────────────────────────────────────────────
\subsection{ข้อมูลพื้นฐานของหลักสูตร}

\subsubsection*{ข้อมูลทั่วไป}

\textbf{ชื่อหลักสูตร (ภาษาไทย)}\\
วิศวกรรมศาสตรมหาบัณฑิต สาขาวิชาวิศวกรรมคอมพิวเตอร์และปัญญาประดิษฐ์

\textbf{ชื่อหลักสูตร (ภาษาอังกฤษ)}\\
Master of Engineering Program in Computer Engineering and Artificial Intelligence

\textbf{ชื่อปริญญา (ภาษาไทย)}\\
วิศวกรรมศาสตรมหาบัณฑิต (วิศวกรรมคอมพิวเตอร์และปัญญาประดิษฐ์)

\textbf{ชื่อปริญญา (ภาษาอังกฤษ)}\\
Master of Engineering (Computer Engineering and Artificial Intelligence)

\textbf{อักษรย่อปริญญา}\\
วศ.ม. (วิศวกรรมคอมพิวเตอร์และปัญญาประดิษฐ์) / M.Eng. (CPE\&AI)

\textbf{คณะที่รับผิดชอบ}\\
คณะเทคโนโลยีอุตสาหกรรม ร่วมกับ คณะวิทยาศาสตร์และเทคโนโลยี

\textbf{สถานภาพหลักสูตร}\\
หลักสูตรใหม่ พ.ศ. 2568

\textbf{สภามหาวิทยาลัยเห็นชอบหลักสูตร}\\
ในคราวประชุมครั้งที่ 4/2568 เมื่อวันที่ 4 เมษายน พ.ศ. 2568

\textbf{จำนวนหน่วยกิตตลอดหลักสูตร}\\
ไม่น้อยกว่า 36 หน่วยกิต

\textbf{รูปแบบหลักสูตร}\\
ปริญญาโท 2 แผน (แผนที่~1 (วิทยานิพนธ์) และแผนที่~2 (สารนิพนธ์))

\textbf{ภาษาที่ใช้จัดการเรียนการสอน}\\
ภาษาไทย

\textbf{ปีที่เริ่มใช้หลักสูตร}\\
ปีการศึกษา 2568 (เริ่มรับนักศึกษา ภาคการศึกษาที่ 2/2568)

\textbf{จำนวนนักศึกษา (ปี 2568)}\\
4 คน (ภาคการศึกษาที่~2/2568)


\subsubsection*{ปรัชญาของหลักสูตร}

\textbf{ปรัชญาของหลักสูตร แผน 1 วิชาการ}\\
บูรณาการความรู้ทางด้านคอมพิวเตอร์และปัญญาประดิษฐ์ผ่านกระบวนการวิจัย เพื่อสร้างองค์ความรู้ใหม่ที่สามารถนำไปประยุกต์ในการแก้ไขปัญหาขององค์กร ชุมชน และท้องถิ่นอย่างยั่งยืน โดยคำนึงถึงหลักคุณธรรมและจริยธรรมวิชาชีพในการพัฒนาเทคโนโลยีปัญญาประดิษฐ์

\textbf{ปรัชญาของหลักสูตร แผน 2 วิชาชีพ}\\
บูรณาการความรู้ทางด้านคอมพิวเตอร์และปัญญาประดิษฐ์ เพื่อพัฒนานวัตกรรมที่ตอบสนองความต้องการและแก้ไขปัญหาขององค์กร ชุมชน และท้องถิ่นอย่างเป็นรูปธรรม ภายใต้กรอบคุณธรรมและจริยธรรมวิชาชีพในการประยุกต์ใช้เทคโนโลยีปัญญาประดิษฐ์

\subsubsection*{ความสำคัญของหลักสูตร}

หลักสูตรวิศวกรรมศาสตรมหาบัณฑิต สาขาคอมพิวเตอร์และปัญญาประดิษฐ์มีความสำคัญอย่างยิ่งในการพัฒนาท้องถิ่นและภูมิภาค โดยมุ่งผลิตบุคลากรที่มีความเชี่ยวชาญสูงด้านเทคโนโลยีสารสนเทศและปัญญาประดิษฐ์ เพื่อตอบสนองความต้องการขององค์กร ชุมชน และท้องถิ่น หลักสูตรนี้เน้นการพัฒนามาตรฐานความรู้และประสบการณ์วิชาชีพ เพื่อให้มหาบัณฑิตสามารถประยุกต์ใช้เทคโนโลยีและบูรณาการศาสตร์ด้านปัญญาประดิษฐ์ให้สอดคล้องกับบริบทขององค์กร ชุมชน และท้องถิ่น ส่งเสริมการวิจัยและพัฒนานวัตกรรมที่แก้ปัญหาและยกระดับคุณภาพชีวิตของคนในชุมชน สนับสนุนการเปลี่ยนผ่านสู่ยุคดิจิทัลของหน่วยงานและองค์กรในท้องถิ่น สร้างผู้นำด้านเทคโนโลยีที่มีความเข้าใจในบริบทท้องถิ่น และเชื่อมโยงองค์ความรู้จากมหาวิทยาลัยสู่การประยุกต์ใช้จริงในพื้นที่ นอกจากนี้ยังปลูกฝังคุณธรรม จริยธรรม และจรรยาบรรณทางวิชาการและวิชาชีพ เพื่อให้มหาบัณฑิตเป็นพลเมืองที่ดี สามารถพัฒนาศักยภาพทางเทคโนโลยีของท้องถิ่นให้สอดคล้องกับบริบทของสังคมในยุคศตวรรษที่ 21 และความต้องการของชุมชนไทย

\subsubsection*{ระบบการจัดการศึกษา}

\textbf{ระบบ}\\
ใช้ระบบไตรภาค (Trimester System) โดย 1 ปีการศึกษาแบ่งออกเป็น 3 ภาคการศึกษาปกติ ได้แก่ ภาคการศึกษาที่ 1, ภาคการศึกษาที่ 2 และภาคการศึกษาที่ 3 โดยแต่ละภาคการศึกษาปกติมีระยะเวลาการศึกษาไม่น้อยกว่า 12 สัปดาห์ ทั้งนี้เป็นไปตามข้อบังคับมหาวิทยาลัยราชภัฏอุตรดิตถ์ ว่าด้วยการศึกษาระดับบัณฑิตศึกษา พ.ศ.~2566 (\hyperlink{ref:AUNQA-4-3-1}{AUNQA-4-3-1})

\textbf{การจัดการศึกษาภาคฤดูร้อน}\\
ไม่มี (เนื่องจากเป็นการจัดการเรียนการสอนในระบบไตรภาคซึ่งแบ่งการศึกษาเป็น 3 ภาคการศึกษาปกติอยู่แล้ว)

\textbf{การเทียบเคียงหน่วยกิตในระบบไตรภาค}\\
ไม่มีการเทียบเคียง

\subsubsection*{การดำเนินการหลักสูตร}

\textbf{วัน-เวลาในการดำเนินการเรียนการสอน}\\
จัดการเรียนการสอนในรูปแบบ \textbf{Weekend Class} (วันเสาร์--อาทิตย์) เวลา 08.30--16.30 น. ตามผลสำรวจความต้องการของผู้มีส่วนได้ส่วนเสีย (83.9\% ต้องการเรียนวันเสาร์--อาทิตย์ --- อ้างอิง \hyperlink{ref:AUNQA-1-4b}{AUNQA-1-4b}) เพื่อรองรับผู้เรียนที่เป็นบุคลากรในภาคอุตสาหกรรม สถานประกอบการ ครู วิศวกร และผู้ปฏิบัติงานในองค์กรต่างๆ

\textbf{คุณสมบัติของผู้เข้าศึกษา}

\begin{enumerate}
  \item สำเร็จการศึกษาระดับปริญญาตรีในสาขาวิศวกรรมคอมพิวเตอร์ วิทยาการคอมพิวเตอร์ เทคโนโลยีสารสนเทศ หรือสาขาที่เกี่ยวข้อง จากสถาบันที่ ก.พ. หรือ สกอ./สป.อว. รับรอง
  \item มีคะแนนเฉลี่ยสะสมไม่ต่ำกว่า 2.50 (กรณีคะแนนเฉลี่ยต่ำกว่า 2.50 ให้พิจารณาประสบการณ์การทำงานที่เกี่ยวข้อง $\geq$2 ปี)
  \item ผ่านการสัมภาษณ์โดยคณะกรรมการที่หลักสูตรแต่งตั้ง
  \item คุณสมบัติอื่นๆ ตามข้อบังคับมหาวิทยาลัยราชภัฏอุตรดิตถ์ ว่าด้วยการศึกษาระดับบัณฑิตศึกษา พ.ศ.~2566
\end{enumerate}

\textbf{ปัญหาของนักศึกษาแรกเข้าและกลยุทธ์การแก้ไข}

\begingroup
\footnotesize
\setlength{\tabcolsep}{4pt}
\renewcommand{\arraystretch}{1.3}
\begin{tabularx}{\linewidth}{|>{\raggedright\arraybackslash}p{4cm}|>{\raggedright\arraybackslash}X|>{\raggedright\arraybackslash}X|}
  \hline
  \rowcolor{gray!15}
  \textbf{ปัญหาของนักศึกษาแรกเข้า} & \textbf{กลยุทธ์ในการดำเนินการแก้ไขปัญหา} & \textbf{ตัวชี้วัดความสำเร็จ} \\
  \hline
  ความรู้ด้านภาษาต่างประเทศ & ให้นักศึกษาลงทะเบียนวิชาภาษาอังกฤษระดับบัณฑิตศึกษา เรียนแบบไม่นับหน่วยกิต & นักศึกษาสอบผ่านรายวิชาภาษาอังกฤษ ซึ่งเป็นไปตามประกาศมหาวิทยาลัยราชภัฏอุตรดิตถ์ เรื่อง เกณฑ์มาตรฐานความสามารถด้านภาษาอังกฤษ สำหรับนักศึกษาระดับบัณฑิตศึกษา พ.ศ.~2561 \\
  \hline
  ความรู้เกี่ยวกับการสืบค้นข้อมูลในการทำวิจัย วิทยานิพนธ์ การอ้างอิงผลการวิจัยที่เกี่ยวข้อง & ฝึกให้นักศึกษาเข้าไปสืบค้นข้อมูลที่เกี่ยวข้องกับงานวิจัย ที่เกี่ยวข้องกับหัวข้อในการทำวิจัย วิทยานิพนธ์ & นักศึกษาเข้าใจและสามารถสืบค้นข้อมูลสำหรับอ้างอิงในการทำวิจัย วิทยานิพนธ์ได้ \\
  \hline
\end{tabularx}
\endgroup

\textbf{แผนการรับนักศึกษาและผู้สำเร็จการศึกษาในระยะ 5 ปี}

\begingroup
\footnotesize
\setlength{\tabcolsep}{4pt}
\renewcommand{\arraystretch}{1.2}
\begin{tabularx}{\linewidth}{|>{\raggedright\arraybackslash}X|c|c|c|c|c|}
  \hline
  \rowcolor{gray!15}
  \textbf{จำนวนนักศึกษา} & \textbf{2568} & \textbf{2569} & \textbf{2570} & \textbf{2571} & \textbf{2572} \\
  \hline
  รับเข้า (เป้าหมาย) & 10 & 10 & 10 & 10 & 10 \\
  \hline
  รับเข้า (จริง) & \textbf{4} & --- & --- & --- & --- \\
  \hline
  คาดว่าจะสำเร็จการศึกษา & --- & --- & 10 & 20 & 30 \\
  \hline
\end{tabularx}
\endgroup

\textbf{ระบบการศึกษา}\\
ระบบการศึกษาเป็นแบบชั้นเรียน (Class-based) จัดในวันเสาร์--อาทิตย์ พร้อมการสนับสนุนผ่านระบบ Online Learning (Microsoft Teams + Google Classroom) เพื่อให้นักศึกษาเข้าถึงเนื้อหาและทรัพยากรการเรียนรู้นอกเวลาเรียนได้

\textbf{การเทียบโอนหน่วยกิต รายวิชา และการลงทะเบียนข้ามมหาวิทยาลัย}\\
เป็นไปตามข้อบังคับมหาวิทยาลัยราชภัฏอุตรดิตถ์ ว่าด้วยการศึกษาระดับบัณฑิตศึกษา พ.ศ.~2566 (\hyperlink{ref:AUNQA-4-3-1}{AUNQA-4-3-1}) ในกรณีที่นักศึกษาเคยศึกษารายวิชาในระดับบัณฑิตศึกษาจากสถาบันอื่นที่ สป.อว. รับรอง สามารถยื่นเรื่องเทียบโอนได้ ตามที่คณะกรรมการบริหารหลักสูตรพิจารณา

\subsubsection*{ผู้มีส่วนได้ส่วนเสียของหลักสูตร}

หลักสูตรได้สำรวจความต้องการและความคาดหวังจากผู้มีส่วนได้ส่วนเสียโดยใช้แบบสอบถาม ซึ่งมีผู้ตอบทั้งสิ้น \textbf{31 คน} ประกอบด้วย ศิษย์เก่า นักศึกษาปัจจุบัน อาจารย์ผู้สอนในหลักสูตร และหน่วยงานภาครัฐและเอกชน อาทิ

\begin{itemize}
  \item โรงพยาบาลอุตรดิตถ์
  \item โรงพยาบาลคีรีมาศ
  \item สำนักงานสาธารณสุขจังหวัดอุตรดิตถ์
  \item สำนักงานเกษตรและสหกรณ์จังหวัดอุตรดิตถ์
  \item บริษัท ศักดิ์สยามลิสซิ่ง จำกัด
  \item บริษัท จัดหางาน จ๊อบบีเคเค ดอท คอม จำกัด
  \item มหาวิทยาลัยราชภัฏอุตรดิตถ์
  \item มหาวิทยาลัยราชภัฏพิบูลสงคราม
\end{itemize}

ข้อมูลจากการสำรวจนำมาวิเคราะห์ร่วมกับนโยบายและยุทธศาสตร์ อววน. พ.ศ. 2563--2570 และแผนยุทธศาสตร์มหาวิทยาลัยราชภัฏอุตรดิตถ์ 2565--2569 ฉบับปรับปรุง เพื่อใช้ออกแบบหลักสูตรทั้งแผน 1 (แบบวิชาการ --- เน้นวิจัย ตอบสนองหน่วยงานภาครัฐและการสาธารณสุข) และแผน 2 (แบบวิชาชีพ --- เน้นการประยุกต์ใช้ ตอบสนองหน่วยงานภาคเอกชน)

\subsubsection*{กลุ่มคู่ความร่วมมือ}

ตามที่ปรากฏใน มคอ.2 (P100--P101) ระบุว่า ``รูปแบบและความร่วมมือการจัดการศึกษาฝึกประสบการณ์วิชาชีพ: ไม่มี'' เนื่องจากเป็นหลักสูตรระดับบัณฑิตศึกษาที่เน้นการทำวิทยานิพนธ์/สารนิพนธ์ในบริบทขององค์กรที่นักศึกษาทำงานอยู่ อย่างไรก็ตาม หลักสูตรได้สร้างความร่วมมือเบื้องต้นกับ \textbf{หน่วยงานภาครัฐและเอกชนในพื้นที่} ที่เข้าร่วมการสำรวจ Stakeholder Needs จำนวน~31~คน (อ้างอิงรายชื่อในหัวข้อ ``ผู้มีส่วนได้ส่วนเสีย'') เพื่อรับฟังความต้องการในการพัฒนาหัวข้อวิทยานิพนธ์/สารนิพนธ์ที่สอดคล้องกับโจทย์จริงในพื้นที่ ทั้งนี้เอกสารการลงนามความร่วมมือ (MOU) ระหว่างมหาวิทยาลัยฯ กับสถานประกอบการปรากฏใน \textbf{ภาคผนวก จ ของ มคอ.2}

\subsubsection*{อาชีพหลังสำเร็จการศึกษา}

มหาบัณฑิตสามารถประกอบอาชีพได้ดังนี้

\begin{enumerate}
  \item วิศวกรปัญญาประดิษฐ์ (Artificial Intelligence Engineer)
  \item วิศวกรการเรียนรู้ของเครื่อง (Machine Learning Engineer)
  \item นักวิจัยด้านปัญญาประดิษฐ์ (Artificial Intelligence Researcher)
  \item วิศวกรประมวลผลภาษาธรรมชาติ (Natural Language Processing Engineer)
  \item วิศวกรระบบวิเคราะห์ข้อมูลอัจฉริยะ (Intelligent Data Analytics Engineer)
  \item วิศวกรระบบคอมพิวเตอร์ (Computer Systems Engineer)
  \item วิศวกรเครือข่าย (Network Engineer)
  \item วิศวกรความปลอดภัยทางไซเบอร์ (Cybersecurity Engineer)
  \item วิศวกรซอฟต์แวร์ฝังตัว (Embedded Software Engineer)
  \item วิศวกรคลาวด์คอมพิวติ้ง (Cloud Computing Engineer)
\end{enumerate}

\subsubsection*{วัตถุประสงค์ของหลักสูตร (PEOs)}

\noindent\textbf{แผน 1 แบบวิชาการ}
\begin{enumerate}
  \item PEO1 มีทักษะในด้านการใช้งานซอฟต์แวร์หรือฮาร์ดแวร์ที่เหมาะสม ในด้านการออกแบบและพัฒนาองค์ความรู้
  \item PEO2 มีความสามารถในการพัฒนาองค์ความรู้เพื่อแก้ปัญหาในองค์กร โดยใช้ความรู้ทางวิศวกรรมคอมพิวเตอร์และปัญญาประดิษฐ์
  \item PEO3 มีคุณธรรมและจรรยาบรรณวิชาชีพทางวิชาการ
  \item PEO4 มีทักษะในการพัฒนาตนเองและการเรียนรู้ตลอดชีวิต
\end{enumerate}

\noindent\textbf{แผน 2 แบบวิชาชีพ}
\begin{enumerate}
  \item PEO1 มีทักษะในด้านการใช้งานซอฟต์แวร์หรือฮาร์ดแวร์ที่เหมาะสม ในด้านการออกแบบและพัฒนาองค์ความรู้
  \item PEO2 มีความสามารถพัฒนานวัตกรรมเพื่อแก้ปัญหาในองค์กร โดยใช้ความรู้ทางวิศวกรรมคอมพิวเตอร์และปัญญาประดิษฐ์
  \item PEO3 มีคุณธรรมและจรรยาบรรณวิชาชีพทางวิชาการ
  \item PEO4 มีทักษะในการพัฒนาตนเองและการเรียนรู้ตลอดชีวิต
\end{enumerate}

\subsubsection*{ผลลัพธ์การเรียนรู้ที่คาดหวัง (Program Learning Outcomes: PLOs)}

\noindent\textbf{แผน 1 แบบวิชาการ}
\begin{itemize}
  \item \textbf{PLO1} ใช้ซอฟต์แวร์หรือฮาร์ดแวร์ทางคอมพิวเตอร์ในการสร้างสรรค์ผลงาน องค์ความรู้ และนวัตกรรม
  \item \textbf{PLO2} ประยุกต์ใช้ความรู้ในสาขาวิชาวิศวกรรมคอมพิวเตอร์เพื่อแก้ไขปัญหาความต้องการขององค์กรชุมชน และท้องถิ่น
  \item \textbf{PLO3} สร้างองค์ความรู้ทางด้านวิศวกรรมคอมพิวเตอร์และปัญญาประดิษฐ์ที่ตอบสนองความต้องการขององค์กร
  \item \textbf{PLO4} ปฏิบัติตามหลักจรรยาบรรณวิชาชีพวิศวกรรมคอมพิวเตอร์
  \item \textbf{PLO5} สามารถคิดวิเคราะห์ การสื่อสาร การทำงานร่วมกัน การปรับตัว และการเรียนรู้ตลอดชีวิต
\end{itemize}

\noindent\textbf{แผน 2 แบบวิชาชีพ}
\begin{itemize}
  \item \textbf{PLO1} ใช้ซอฟต์แวร์หรือฮาร์ดแวร์ทางคอมพิวเตอร์ในการสร้างสรรค์ผลงาน องค์ความรู้ และนวัตกรรม
  \item \textbf{PLO2} ประยุกต์ใช้ความรู้ในสาขาวิชาวิศวกรรมคอมพิวเตอร์เพื่อแก้ไขปัญหาความต้องการขององค์กรชุมชน และท้องถิ่น
  \item \textbf{PLO3} สร้างต้นแบบนวัตกรรมด้านวิศวกรรมคอมพิวเตอร์และปัญญาประดิษฐ์ที่ตอบสนองความต้องการขององค์กร
  \item \textbf{PLO4} ปฏิบัติตามหลักจรรยาบรรณวิชาชีพวิศวกรรมคอมพิวเตอร์
  \item \textbf{PLO5} สามารถคิดวิเคราะห์ การสื่อสาร การทำงานร่วมกัน การปรับตัว และการเรียนรู้ตลอดชีวิต
\end{itemize}

\subsubsection*{โครงสร้างของหลักสูตร}

\textbf{จำนวนหน่วยกิตรวมตลอดหลักสูตร: ไม่น้อยกว่า 36 หน่วยกิต}

\vspace{4pt}
\begingroup
\footnotesize
\setlength{\tabcolsep}{4pt}
\renewcommand{\arraystretch}{1.2}
\begin{tabularx}{\linewidth}{l>{\raggedright\arraybackslash}X r l r l}
  \multicolumn{6}{l}{\textbf{แผนการศึกษา 1 (แบบวิชาการ --- วิทยานิพนธ์)}} \\
  \multicolumn{6}{l}{} \\
  \textbf{(1)} & หมวดวิชาเฉพาะด้านบังคับ เรียนไม่น้อยกว่า & & & \textbf{18} & หน่วยกิต \\
  & \quad (1.1) วิชาเฉพาะด้านบังคับ 6 รายวิชา & & & 18 & หน่วยกิต \\
  \textbf{(2)} & หมวดวิชาเฉพาะด้านเลือก เรียนไม่น้อยกว่า & & & \textbf{6} & หน่วยกิต \\
  \textbf{(3)} & หมวดวิทยานิพนธ์ & & & \textbf{12} & หน่วยกิต \\
  & \quad (3.1) วิทยานิพนธ์ 1 (เสนอหัวข้อ) & & & 3 & หน่วยกิต \\
  & \quad (3.2) วิทยานิพนธ์ 2 (เก็บข้อมูล) & & & 3 & หน่วยกิต \\
  & \quad (3.3) วิทยานิพนธ์ 3 (วิทยานิพนธ์ฉบับสมบูรณ์) & & & 6 & หน่วยกิต \\
  \textbf{(4)} & หมวดวิชาเสริม (ไม่นับหน่วยกิต) & & & (6) & หน่วยกิต \\
  & \quad (4.1) ภาษาอังกฤษสำหรับบัณฑิตศึกษา & & & (3) & หน่วยกิต \\
  & \quad (4.2) การเขียนวิทยานิพนธ์หรือสารนิพนธ์ & & & (3) & หน่วยกิต \\
  \multicolumn{4}{r}{\textbf{รวมหน่วยกิตแผน 1 ไม่น้อยกว่า}} & \textbf{36} & \textbf{หน่วยกิต} \\
  \multicolumn{6}{l}{} \\
  \multicolumn{6}{l}{\textbf{แผนการศึกษา 2 (แบบวิชาชีพ --- สารนิพนธ์)}} \\
  \multicolumn{6}{l}{} \\
  \textbf{(1)} & หมวดวิชาเฉพาะด้านบังคับ เรียนไม่น้อยกว่า & & & \textbf{18} & หน่วยกิต \\
  & \quad (1.1) วิชาเฉพาะด้านบังคับ 6 รายวิชา & & & 18 & หน่วยกิต \\
  \textbf{(2)} & หมวดวิชาเฉพาะด้านเลือก เรียนไม่น้อยกว่า & & & \textbf{12} & หน่วยกิต \\
  \textbf{(3)} & หมวดสารนิพนธ์ & & & \textbf{6} & หน่วยกิต \\
  & \quad (3.1) สารนิพนธ์ 1 (เสนอหัวข้อ) & & & 3 & หน่วยกิต \\
  & \quad (3.2) สารนิพนธ์ 2 (สารนิพนธ์ฉบับสมบูรณ์) & & & 3 & หน่วยกิต \\
  \textbf{(4)} & หมวดวิชาเสริม (ไม่นับหน่วยกิต) & & & (6) & หน่วยกิต \\
  \multicolumn{4}{r}{\textbf{รวมหน่วยกิตแผน 2 ไม่น้อยกว่า}} & \textbf{36} & \textbf{หน่วยกิต} \\
\end{tabularx}
\endgroup

\vspace{6pt}
\subsubsection*{จำนวนนักศึกษาปัจจุบันในแต่ละชั้นปี ปีการศึกษา 2568}

\begingroup
\footnotesize
\setlength{\tabcolsep}{4pt}
\renewcommand{\arraystretch}{1.3}
\begin{tabularx}{\linewidth}{|>{\raggedright\arraybackslash}X|c|c|c|c|c|}
  \hline
  \rowcolor{gray!15}
  \multirow{2}{*}{\textbf{แผนการศึกษา}} & \multicolumn{4}{c|}{\textbf{ระดับชั้นปี (ปีที่รับเข้า)}} & \multirow{2}{*}{\textbf{รวม (คน)}} \\
  \cline{2-5}
  & \textbf{ปี 1 (2568)} & \textbf{ปี 2 (2567)} & \textbf{ปี 3 (2566)} & \textbf{ตกค้าง} & \\
  \hline
  แผน 1 (วิทยานิพนธ์) & 0 & --- & --- & --- & 0 \\
  \hline
  แผน 2 (สารนิพนธ์) & 4 & --- & --- & --- & 4 \\
  \hline
  \textbf{รวม} & \textbf{4} & \textbf{---} & \textbf{---} & \textbf{---} & \textbf{4} \\
  \hline
\end{tabularx}
\endgroup

\vspace{4pt}
\footnotesize \textit{หมายเหตุ หลักสูตรเป็นหลักสูตรใหม่ พ.ศ.~2568 เริ่มรับนักศึกษารุ่นแรกในภาคการศึกษาที่ 2/2568 จำนวน 4 คน ทั้งหมดเลือกแผนการศึกษาที่ 2 (สารนิพนธ์); ยังไม่มีนักศึกษาในแผน 1 และยังไม่มีนักศึกษาตกค้าง}

\vspace{6pt}
\textbf{กลุ่มผู้เรียน} นักศึกษาปัจจุบันที่ต้องการนำความรู้ไปใช้ในการทำงานและพัฒนาอาชีพ ศิษย์เก่าและผู้สำเร็จการศึกษาระดับปริญญาตรีในสาขาที่เกี่ยวข้อง รวมถึงบุคลากรของหน่วยงานภาครัฐและเอกชนในจังหวัดอุตรดิตถ์และพื้นที่ใกล้เคียง (อ้างอิง มคอ.2 P144)
